\begin{abstract}
Reconstructing the temporal sequence of mutations and genomic events during tumour expansion is a fundamental challenge in cancer data science. In this dissertation, a Bayesian model is developed and validated for the joint inference of evolutionary timings and clonal growth rates from single-cell sequencing data. The statistical inference approach uses four main factors: the number of somatic mutations accumulated along each cell clone, the occurrence of copy-number alterations (large-scale genomic duplications), the abundance of distinct cell populations, and the topology of the phylogenetic tree relating them. A hierarchical prior structure is adopted to share statistical strength across closely related clones while preserving variability across distinct subpopulations. Somatic mutations are modelled as a stochastic counting process, with intensity depending on the background mutation rate, the length of the genomic region, an event-specific coefficient, the clone-specific growth rate, and the time of evolution of the clone. The timing of copy-number alterations is inferred by partitioning the mutations mapping to each affected genomic region into those arising before and after the event, exploiting the discontinuity in the mutation accumulation process induced by the copy-number change as a temporal signal.

The inference framework is assessed through simulation studies of increasing complexity. Under a simple generative scenario the model recovers ground-truth timings and growth rates with good accuracy. A more demanding evaluation is carried out using synthetic data from ProCESS, a cancer evolution simulator, which allows examination of how deviations from the modelling assumptions propagate into the inference outcomes.

\end{abstract}